\section{Discussion}

\subsection{The Symmetric Gradient}

The authoring ablation produces the same non-monotonic gradient shape as the reviewer ablation in Paper~I: more context does not monotonically improve quality, and the largest single inflection occurs not at the introduction of domain data but at the introduction of a design philosophy. The A0-to-A3 climb (+9.3 points across four conditions) mirrors the C0-to-C3 climb in Paper~I almost exactly in proportional terms; the plateau-and-retreat at A4–A6 mirrors the C4–C6 flattening on the reviewer side.

The A2 result is the clearest structural parallel to Paper~I's C2 finding. In the reviewer study, principles named without operational definitions produced critiques that were measurably worse-calibrated than richer contexts. Here, domain data without a design framework produces puzzles that score only 0.6 points above bare task specification (28.3 vs.\ 27.7) despite extensive world-data citations. In both cases, providing raw content — factual data in A2, principle names in C2 — without a conceptual framework for how to use it yields output less effective than a modest increase in structural guidance would produce. The mechanism is the same in both directions: domain expertise loaded into a prompt does not automatically become domain expertise in the output. The profile's framework determines what the model does with the information it holds.

The A4$\approx$A5 result is an equally precise parallel to Paper~I's C4$\approx$C3 finding, which showed that a reviewer's worldview was evaluatively as rich as a full operational checklist. For authoring, a designer's generative starting position — where they begin, what question they ask first, what the finish condition looks like — captures the overwhelming majority of the quality value of a complete construction profile. The 1.3-point gap between A4 and A5 reflects a specific and narrow contribution: verification rigor. The full profile's construction log eliminates ambiguous item candidates and cross-checks competing answers, which protects execution quality within a chosen mechanism. It does not elevate the mechanism, alter the design strategy, or produce qualitatively different puzzles. For a designer with strong domain instincts, worldview is the rate-limiting resource; the checklist is insurance.

The symmetry between the two ablations has a broader implication. The mechanisms that determine profile quality — framework versus content, generative worldview versus procedural list — are not specific to the evaluation task. They appear to be structural properties of how profiles engage model behavior, applicable whenever a profile must guide a complex, open-ended generation task. The puzzle hunt domain provides unusually clean measurement conditions, but the gradient shape is likely general.

\subsection{The Lens-Type Transfer Model}

The three-designer comparison yields a \emph{proposed} three-category taxonomy of how reviewing expertise transfers to authoring that is determined by lens type rather than by the designer's overall profile quality or depth. With one designer per category, this taxonomy cannot be empirically validated from the present data; it is offered as a theoretically motivated classification scheme requiring replication with a broader designer sample. This result is nonetheless theoretically significant: it means that the question ``can this reviewer profile be repurposed for authoring?'' has a principled answer based on the profile's internal structure, not its surface sophistication.

\textit{Operational lenses} (Katz) are inherently generative. Each named test in Katz's reviewing vocabulary has a direct authoring-time analog that is neither a metaphor nor an approximation. ``Would I want to solve this?'' inverts to a constraint: design something you yourself would want to solve, and evaluate candidates by applying that test during construction. ``Can the meta be short-circuited?'' becomes: ensure no subset of four items resolves the answer without requiring the fifth. The Computer Test — eliminating clues a search engine can answer without domain knowledge — becomes a candidate filter applied before any clue is written. These inversions are lossless: the reviewing test and the authoring constraint are logically equivalent, distinguished only by the direction of application. The A6-Katz result (35.3/45, the WOLOLO puzzle) confirms this: the reviewer-as-author exercise produced the strongest puzzle among all Katz conditions precisely because the reviewing discipline sharpened item selection rather than being abandoned or approximated.

\textit{Construction lenses} (Snyder) are already authoring-compatible. Snyder's reviewing questions — is each element load-bearing? does the solve path hold in exactly one direction? are there cross-clue dependencies that reward deductive rather than guessing strategy? — are indistinguishable from authoring instructions. The A5-to-A6 gap for Snyder is $-0.7$ points, within noise, because both profiles execute the same underlying checklist. For construction-oriented experts, the authoring/reviewing distinction is not a cognitive distinction; they apply the same tests in both directions. This collapses the practical difference between the A5 and A6 profiles for this lens type.

\textit{Perceptual lenses} (Young) are terminal and format-dependent. Young's reviewing questions — does the visual hierarchy communicate structure before the content does? is the layout doing independent semantic work? does the visual grammar of the puzzle reinforce its domain? — require finished material in a format that can carry visual and experiential information. In a text-only acrostic, these checks have no application surface. The A5-Young puzzle (34.3/45) is strong but proceeds without Young's primary evaluative lens engaging at all; she substitutes phenomenological clue texture for her usual perceptual criteria. The A6-Young result (38.0/45, the highest score in the dataset) is therefore the finding that requires the most careful interpretation: it is not evidence that perceptual lenses transfer. It is evidence that blocking a designer's default approach — by forcing the reviewer-as-author exercise to surface the inapplicability of visual checks — can produce a productive substitution. Young found a format-compatible analog: second-person present-tense narrative that places the solver inside the AoE2 world rather than outside it. The perceptual instinct (make the solver feel inside the domain) found expression through narrative grammar because visual grammar was unavailable.

The practical implication for profile construction follows directly. When building authoring profiles from reviewing profiles: operational tests should be preserved and inverted; construction tests are transferable as-is; perceptual checks should be replaced by their generative precursors — not ``does the layout communicate structure?'' but ``design the layout so its structure communicates before the content does.'' The last transformation requires identifying what the perceptual check is ultimately verifying (domain presence, experiential immersion, formal coherence) and constructing a generative instruction that can be satisfied during creation rather than only assessed after the fact.

A profile's lens type can be tentatively classified from its review lens section alone, without observing authoring behavior: count the proportion of questions framed as named operations with defined inputs and outputs (operational), as uniqueness or load-bearing tests (construction), and as perceptual or aesthetic assessments requiring finished material (perceptual). This procedure allows classification prior to any authoring study.

Two alternative explanations for the A5-to-A6 gap differences deserve acknowledgment. First, the gaps may reflect not lens type but profile length: Katz's C8 reviewing profile is longer than his authoring profile, which could independently account for A6-Katz outperforming A5-Katz. Second, the Young paradox (A6 > A5 despite an inapplicable lens) may reflect creative constraint effects specific to the text-acrostic format, and may not generalize to formats where her visual lens can be applied directly. Controlling for profile length and format compatibility are priorities for future work.

\subsection{The Format Ceiling and the Riven Standard}

All A-series puzzles in the authoring ablation score Riven Standard $\leq 3$, regardless of profile depth or construction rigor. The first-letter acrostic mechanism is imported from puzzle hunt tradition and is not native to AoE2's game logic: the extraction step (take the first letter of each item name) has no analog in how AoE2 structures knowledge, rewards expertise, or produces meaning. The double-weighting of Elegance and Reading Reward in the $/45$ rubric penalizes format-imported mechanisms by design, since both dimensions reward domain-mechanism integration — the sense that the puzzle could not have been designed in any other domain. The Riven Standard rubric dimension makes this penalty explicit: a score of 3 indicates that the domain is content (the items are real AoE2 items) but not mechanism (the extraction logic is domain-agnostic).

The A6-Katz and A6-Young puzzles score Riven Standard 4, and A6-Young scores 5 — the only Riven 5 in the dataset. The distinction is mechanism adaptation: A6-Young constructs a coherent tactical scenario in which each item plays a role in a specific defensive situation, so the answer \textsc{tower} names not just a structure but the strategic position the solver has been defending across all five clues. This is the closest any evaluated puzzle comes to mechanism-domain integration without abandoning the acrostic format entirely.

The ceiling finding has a counterintuitive implication for profile depth. Profile depth raises the floor of quality within a fixed mechanism: the difference between A0 (21.7/45) and A5 (32.0/45) is real and large. But mechanism choice determines the ceiling, and no authoring profile in the A-series — not even the most sophisticated — can earn a high Riven Standard score when the extraction mechanism is format-imported. The companion persona study demonstrates this directly: a minimal two-word authoring persona — ``you are an artist'' — outperforms all A-series conditions (40.3/45, Riven Standard 4.3/5) by generating a mechanism native to the game's visual and archaeological register (damaged inscriptions with one missing letter per item, each gap corresponding to a domain-specific property). The artist persona did not produce a better-executed puzzle on any individual dimension; it produced a different class of mechanism. Profile depth and mechanism novelty are orthogonal contributions to quality, and the rubric is sensitive to both.

\subsection{Implications for Computational Creativity}

The symmetric gradient — authoring mirrors reviewing in both shape and inflection points — speaks directly to the \emph{evaluation problem} in computational creativity \cite{colton2012computational}: if the same competence profile structure governs both generation quality and evaluation quality, then systems that evaluate well and systems that generate well may require the same underlying profile architecture, differing only in the direction of application. This has design implications: a reviewing profile is not a secondary artifact but a specification of the same latent competence that would govern authoring.

The A4$\approx$A5 finding connects to Boden's account of \emph{exploratory creativity} \cite{boden2004creative}: a designer who has internalized a generative worldview explores the same conceptual space as one equipped with a full operational checklist, because the worldview implicitly encodes the boundaries of that space. The checklist makes boundary conditions explicit; the worldview makes them operational. Either representation produces equivalent exploration of the design space when the designer is sufficiently expert.

The Young A6 result — where a perceptual lens blocked from its default application produced the highest-scoring output in the dataset — connects to the \emph{creative constraint} literature \cite{stokes2006}: constraints on a designer's natural approach can produce more creative outputs by blocking the default path and forcing a less-explored one. Young's perceptual instinct could not operate in the text-only format; the resulting substitution (narrative-grammar immersion for visual-grammar immersion) accessed a region of the design space her default approach would not have reached.

These findings position profile structure as a theoretically tractable variable in CC system design, not merely a prompt-engineering parameter. The companion persona-authoring-traits study takes this further by examining which profile features are predictive of mechanism novelty, the dimension most directly connected to Boden's transformational creativity account.

\subsection{Limitations}

Four limitations constrain the scope of these conclusions. First, all conditions use the same target answer (\textsc{tower}) and the same mechanism type (knowledge extraction via first-letter acrostic). The gradient findings may not generalize to other mechanism types — specifically to mechanisms where the extraction logic is already domain-native, which would lower the format ceiling for all conditions and potentially alter the relative contribution of profile depth. Second, the Young lens finding is format-specific. In a visually structured puzzle format — a grid, a map overlay, a layout-dependent extraction — Young's A6 profile would likely fail to achieve the forced substitution that produced the highest score in this dataset, and we predict it would underperform A5-Young in such formats. Third, the A5 and A6 authoring profiles used in the designer comparison were constructed for this study. They capture each designer's published design philosophy and stated principles but have not been validated against the designers' actual construction practice. The profiles are models of the designers, not the designers themselves. Fourth, evaluation calibration requires same-session scoring; all 13 puzzles in this study were read and ranked before any scores were assigned, following the calibration protocol established in Paper~I. Scores from evaluation sessions conducted under different calibration conditions are not directly comparable to the values reported here, and earlier pilot scores for overlapping conditions have been superseded by the canonical c8-unified results.

A further limitation is that the three C8 panel reviewers (Katz, Snyder, Young) are also the designers whose profiles are used in Group B. This creates a potential self-evaluation confound: each reviewer evaluates puzzles partially authored under their own profile. We expect this produces conservative scores on A5 conditions (familiar constraints) and lenient scores on A6 conditions (recognizing their own lens), but the direction is uncertain.
